\documentclass[11pt]{article}
\usepackage[margin=1in]{geometry}
\usepackage{amsmath,amssymb,amsthm}
\usepackage{graphicx}
\usepackage[hidelinks]{hyperref}

\newtheorem{theorem}{Theorem}
\newtheorem{corollary}[theorem]{Corollary}
\newtheorem{lemma}[theorem]{Lemma}
\newcommand{\dist}{\operatorname{dist}}
\newcommand{\id}{\operatorname{id}}

\title{Sharp affine Hyers--Ulam constants\\for $\ell_\infty$-products}
\author{A substantial partial result for arXiv:math/0201098}
\date{11 August 2026}

\begin{document}
\maketitle

\begin{center}
\fbox{\parbox{0.91\textwidth}{\textbf{Status: substantial partial result,
likely valid.} Exact sharpness propagates through arbitrary max products.
This settles many new non-Hilbert cases but not the source question for every
Banach space.}}
\end{center}

\section{Source question and notation}

A map $f:E\to E$ is an $\varepsilon$-\emph{nearisometry} if
\[
 \big|\,\|f(x)-f(y)\|-\|x-y\|\,\big|\leq\varepsilon
 \qquad(x,y\in E).
\]
Let $H(E)$ be the infimum of the constants $c$ such that every surjective
$\varepsilon$-nearisometry of the real Banach space $E$ is within
$c\varepsilon$ of a surjective affine isometry.  Let $J(E)$ be the Jung
constant: the infimum of $r$ such that every subset of $E$ of diameter at
most $2$ is contained in a ball of radius $r$.

V\"ais\"al\"a's survey \cite{Vaisala} recalls the general upper bound
$H(E)\leq J(E)$ and asks on printed page 4 whether it is sharp for every
Banach space.

\begin{figure}[ht]
\centering
\includegraphics[width=0.94\textwidth]{figures/source_question.png}
\caption{The exact sharpness question in \cite[p.~4]{Vaisala}.}
\end{figure}

Huuskonen and V\"ais\"al\"a \cite{HV} prove $H(E)=J(E)$ for every Hilbert
space and record the same equality for $c$, $c_0$, spaces with $J(E)=1$,
and one regular-hexagon plane.  The following theorem propagates any such
equality through max products.

\section{Idea of the proof}

The max norm decouples both sides of the problem.  Coordinate projections
turn a diameter-two set in a product into diameter-two sets in every factor;
coordinatewise Jung centers then assemble into one bounded product center.
For the lower Hyers--Ulam bound, insert a bad nearisometry in one coordinate
and leave all other coordinates fixed.  The max operation preserves the same
additive error.

The only possible complication is that an ambient isometry might mix
coordinates.  It cannot do so when it approximates the lifted map: after a
standard normalization the lifted map is a bounded distance from the
identity, and every affine isometry a bounded distance from the identity has
identity linear part.  Thus any approximant is merely a translation, whose
bad coordinate recovers the original factor obstruction.

\section{Product theorem}

For a nonempty family $(X_i)_{i\in I}$ of nonzero real Banach spaces, write
\[
 X=\left(\prod_{i\in I}X_i\right)_\infty
 =\left\{(x_i):\sup_i\|x_i\|<\infty\right\},
 \qquad \|(x_i)\|=\sup_i\|x_i\|.
\]

\begin{theorem}[max-product propagation]\label{thm:product}
For $X=(\prod_iX_i)_\infty$,
\[
 J(X)=\sup_{i\in I}J(X_i)
 \quad\text{and}\quad
 H(X)\geq\sup_{i\in I}H(X_i).
\]
Consequently, if
\[
 \sup_iH(X_i)=\sup_iJ(X_i),
\]
then $H(X)=J(X)$.  In particular this conclusion holds whenever
$H(X_i)=J(X_i)$ for every $i$.
\end{theorem}

\begin{proof}
Set $J_*:=\sup_iJ(X_i)$.  Each $X_i$ embeds isometrically as a coordinate
subspace of $X$, and projecting a covering ball in $X$ onto that coordinate
shows $J(X)\geq J(X_i)$.  Hence $J(X)\geq J_*$.

Conversely, let $A\subset X$ have diameter at most $2$, fix $a^0\in A$, and
let $\eta>0$.  For every $i$, the coordinate set
$A_i=\{a_i:a\in A\}$ has diameter at most $2$, so there is $c_i\in X_i$ with
\[
 \sup_{a\in A}\|a_i-c_i\|\leq J(X_i)+\eta\leq J_*+\eta.
\]
Moreover,
\[
 \|c_i\|\leq\|c_i-a_i^0\|+\|a_i^0\|
 \leq J_*+\eta+\|a^0\|,
\]
uniformly in $i$.  Thus $c=(c_i)$ belongs to $X$ and
$A\subset B_X(c,J_*+\eta)$.  Letting $\eta\downarrow0$ proves
$J(X)=J_*$.

We next prove $H(X)\geq H(X_k)$ for a fixed $k\in I$.  Let
$c<H(X_k)$.  By the definition of $H(X_k)$ there are
$\varepsilon>0$ and a surjective $\varepsilon$-nearisometry
$f:X_k\to X_k$ whose distance from every surjective affine isometry exceeds
$c\varepsilon$.

Translate $f$ so that $f(0)=0$.  The fundamental nearisometry theorem
\cite[Theorem~2.2]{Vaisala} supplies a surjective linear isometry
$U:X_k\to X_k$ with $d(f,U)\leq2\varepsilon$.  Replacing $f$ by
$g=U^{-1}f$ preserves all distances from the affine-isometry family and gives
\[
 d(g,\id_{X_k})<\infty.
\]

Lift $g$ to $G:X\to X$ by
\[
 (Gx)_k=g(x_k),\qquad (Gx)_i=x_i\quad(i\neq k).
\]
This map is surjective.  If
$a=\|x_k-y_k\|$, $a'=\|g(x_k)-g(y_k)\|$, and
$b=\sup_{i\neq k}\|x_i-y_i\|$, then
\[
 \big|\,\|Gx-Gy\|-\|x-y\|\,\big|
 =|\max\{a',b\}-\max\{a,b\}|
 \leq|a'-a|\leq\varepsilon.
\]
Hence $G$ is an $\varepsilon$-nearisometry, and
$d(G,\id_X)<\infty$.

Let $T:X\to X$ be a surjective affine isometry with $d(T,G)<\infty$.
By Mazur--Ulam, $Tz=Lz+w$ for a surjective linear isometry $L$.  Since
$d(T,\id_X)<\infty$, for every $z\in X$ the vectors
\[
 t(L-I)z+w=T(tz)-tz
\]
remain bounded as $t\to\infty$.  Therefore $L=I$ and $T$ is translation by
$w=(w_i)$.  Restricting to vectors supported in coordinate $k$ gives
\[
 d(T,G)\geq\sup_{u\in X_k}\|u+w_k-g(u)\|>c\varepsilon,
\]
because $u\mapsto u+w_k$ is an affine isometry of $X_k$.  Thus the lifted
map defeats every ambient affine isometry at the same constant, and
$H(X)\geq c$.  Letting $c\uparrow H(X_k)$ and then taking the supremum over
$k$ proves the second assertion.

Finally the general estimate $H(X)\leq J(X)$ from \cite{Vaisala}, together
with the two formulas just proved, gives the equality under the stated
hypothesis.
\end{proof}

\section{Concrete new cases}

\begin{corollary}\label{cor:examples}
The equality $H(X)=J(X)$ holds in each of the following cases.
\begin{enumerate}
\item $X=(\prod_i H_i)_\infty$ for arbitrary real Hilbert spaces $H_i$.
In particular,
\[
 H(X)=J(X)=\sup_i J(H_i).
\]
For finite-dimensional $H_i$ of dimensions $n_i$, the factor constant is
$\sqrt{2n_i/(n_i+1)}$; an infinite-dimensional factor or unbounded finite
dimensions give supremum $\sqrt2$.
\item $X=c_0\oplus_\infty Y$ for every real Banach space $Y$; here
$H(X)=J(X)=2$.
\item $X=\ell_2^2\oplus_\infty\mathbb R$; here
\[
 H(X)=J(X)=\frac{2}{\sqrt3}.
\]
This three-dimensional max product is not a Hilbert space.
\end{enumerate}
\end{corollary}

\begin{proof}
Hilbert factors satisfy $H=J$ by \cite{HV}.  Also $H(c_0)=J(c_0)=2$ by
\cite[Theorem~3.2]{HV}, while every Banach space has Jung constant at most
$2$.  The final formula uses
$J(\ell_2^2)=\sqrt{4/3}=2/\sqrt3$ and $J(\mathbb R)=1$.
Apply Theorem~\ref{thm:product}.
\end{proof}

\section{Verification, novelty, and limitations}

The proof uses only three externally established facts: Mazur--Ulam, the
fundamental $2\varepsilon$ approximation theorem, and the upper bound
$H(E)\leq J(E)$.  The two product arguments were separately checked for
arbitrary (including uncountable) index sets; boundedness of the assembled
Jung center is the only point not automatic coordinatewise.

A bounded novelty search through 11 August 2026 checked the run's exact-id and
keyword indexes; exact web searches for ``Hyers--Ulam constant direct sum'',
``$H(E)=J(E)$ nearisometry product'', and variants; all nine OpenAlex-indexed
papers citing the source survey; and all seven OpenAlex-indexed papers citing
\cite{HV}.  No direct-sum or max-product formula for these constants was
located.  Novelty confidence is moderate rather than definitive because the
topic has sparse and partly pre-arXiv literature.

The full source question remains open.  For finite $p$, the same lifting gives
only
\[
 H(X\oplus_pY)\geq\max\{H(X),H(Y)\},
\]
whereas the Jung constant can be strictly larger: already
$J(\mathbb R\oplus_2\mathbb R)=2/\sqrt3>1$.  General norm-one complemented
subspaces also lack the max metric needed to preserve a uniform additive
error.  A full solution therefore needs a coupled construction realizing
general Jung-extremal configurations, not merely coordinatewise witnesses.

\paragraph{Human-review recommendation.}
Check especially (i) normalization of the factor nearisometry by a linear
isometry, (ii) the assertion that a finite-distance affine isometry has
identity linear part, and (iii) boundedness of the product Jung center.

\begin{thebibliography}{9}
\bibitem{Vaisala}
J.~V\"ais\"al\"a,
\emph{A survey of nearisometries},
in \emph{Papers on Analysis}, Reports of the Department of Mathematics and
Statistics, University of Jyv\"askyl\"a 83 (2001), 305--315;
arXiv:\href{https://arxiv.org/abs/math/0201098}{math/0201098}.

\bibitem{HV}
T.~Huuskonen and J.~V\"ais\"al\"a,
\emph{Hyers--Ulam constants of Hilbert spaces},
Studia Math. 153 (2002), 31--40;
DOI: \href{https://doi.org/10.4064/sm153-1-3}{10.4064/sm153-1-3}.
\end{thebibliography}

\end{document}
